\documentclass[a4paper,12pt]{letter}

\usepackage{amsmath}
\usepackage{amssymb}
\usepackage{nopageno}
\usepackage{amsmath}
\usepackage{enumitem}
% \usepackage{mathtools}
%\usepackage{eqnarray,amsmath}
\usepackage[a4paper, total={7.5in, 11.25in}]{geometry}

\begin{document}

Name: \hrulefill\ Neptun code: \hrulefill

\bigskip

\begin{center}
(KTXFI2EBNF) Physics I. Exam~1 \\
May~26,~2026
\end{center}

\bigskip
 
\begin{enumerate}[label=\Roman*.]
   
\item Mechanics

\begin{enumerate}[label=\arabic*.]
\setcounter{enumii}{0} % Számláló beállítása

\item What does Newton’s first law state? What does Newton’s second law describe? What is Newton’s third law? \hfill (5~points)

% (Lecture~on~2026-02-16, slide~17)

\item What are conservative forces? What are dissipative forces? \hfill (5~points)

%(Lecture~on~2026-02-16, slide 6)

\item What is momentum? What is conservation of momentum? \hfill (5~points)

% (Lecture~on~2026-02-16, slide~15) (Lecture~on~2026-02-16, slide~18)

  
\item What is escape velocity? \hfill (5~points)

%(Lecture~on~2026-02-16, slide~31)
  
\end{enumerate}

% ------------------------------------------------------------------------------------------------------------------------------------

\item Thermodynamics

\begin{enumerate}[label=\arabic*.]
\setcounter{enumii}{4} % Számláló beállítása

\item What does the first law of thermodynamics state? What is the second law of thermodynamics?

\hfill (5~points)

% (2026-03-16, slide 11)  (2026-03-23, slide 11)

\item What is thermal expansion? What is linear expansion? What is volumetric expansion? \hfill (5~points)

%  (2026-03-09, slide 5),  (2026-03-09, slide 5) (2026-03-09, slide 7)


\item What is the ideal gas model? What is the ideal gas equation conceptually? \hfill (5~points)

% (2026-03-09, slide 20)  (2026-03-09, slide 21)

\item What is the Carnot cycle? \hfill (5~points)

% (2026-03-23, slide 9)

\end{enumerate}
  
% ------------------------------------------------------------------------------------------------------------------------------------

\item Electrostatics, Waves, and Optics

\begin{enumerate}[label=\arabic*.]
\setcounter{enumii}{8} % Számláló beállítása
  
\item What is Coulomb’s law? \hfill (5~points)

% (2026-04-13, slide 13)
  
\item What is the Doppler effect? \hfill (5~points)

% (2026-03-02, slide 32)

\end{enumerate}

\end{enumerate}

% ------------------------------------------------------------------------------------------------------------------------------------

Grades: 0–25~points: Fail (1), 26–30 points: Pass (2), 31–37 points: Satisfactory (3), 38–45~points: Good~(4), 46–50 points: Excellent (5).

\rightline{Dr.~Gabor FACSKO}
\rightline{\textit{facsko.gabor@uni-obuda.hu}}

\vfill

\end{document}
